\section{Experimental setup details}
\label{sec:appx-eval}

\subsection{Experimental setup for \Cref{sec:char-obs}}
\label{sec:appx-setup-obs}

\paragraph{Observation 1.}
Routing skew and decode bubbles use Qwen3-30B-A3B on vLLM with
\ac{TP}$=$\ac{EP}$=2$ on two A100 GPUs, eager mode, and no load
balancing. Skew is the most-loaded GPU's routed tokens divided by the
mean, summed over all \ac{MoE} layers, from vLLM's EPLB balancedness
logging with no redundant experts and no rebalancing. Each dataset
sends 30 requests one at a time, so each request contributes one
prefill step and one decode step per generated token. Decode bubbles
are measured from GPU kernel timestamps recorded with
\texttt{torch.profiler}.

\paragraph{Observation 2.}
Load-balancing results (\Cref{fig:lb-sweep}) use DeepSeek-V2-Lite on
SGLang with 8$\times$H100, \ac{EP}$=8$, BF16, and \textsc{arXiv}
prompts. Each point is the median over 1,248 samples: six prefill
forwards, 26 \ac{MoE} layers, and eight GPUs.

\paragraph{Observation 3.}
The colocation systems and the planner experiment are described in
Appendix~\ref{sec:appx-gap}.

\subsection{Evaluation setup}
\label{sec:appx-setup}

This section gives the full configuration summarized in \S\ref{sec:eval-method}.

\paragraph{Hardware and software.}
All experiments run on one node with four NVIDIA A100-SXM4-80GB GPUs
connected by NVLink. Unless stated otherwise, serving uses vLLM with
\ac{TP}=\ac{EP}=4, BF16 weights, eager execution, chunked prefill and
prefix caching disabled, and a 0.85 GPU-memory budget. \sys{} runs one
fine-tuning worker per GPU under CUDA MPS, as described in
\S\ref{sec:design}, with a 50\,$\mu$s chunk budget and an early-close
margin of $m=60\,\mu$s (\S\ref{sec:design-optC}).

\paragraph{Models.}
We evaluate three MoE architectures with different expert
organizations. Qwen3-30B-A3B (48 layers, 128 experts, top-8) uses only
routed experts, while Qwen1.5-MoE-A2.7B (24 layers, 60 routed experts,
top-4) and DeepSeek-V2-Lite (27 layers, 64 routed experts, top-6) also
contain shared experts. Qwen1.5-MoE's
configuration declares an 8k context; for longer prompts we extend its
rotary cache to the served length identically for every system.

\paragraph{Serving workloads and metrics.}
Our controlled workloads use arXiv summarization prompts truncated to
an exact input length between 1k and 32k tokens, each generating 32
output tokens. Short outputs make the workload prefill-dominated, where
bubbles arise, and resemble short-form tasks such as extraction or
classification; \S\ref{sec:eval-e2e} also replays production traces.

For every model and prompt length, we choose three arrival rates:
\emph{high} is just below the point at which serving alone begins
sustained queueing, while \emph{medium} and \emph{low} are fixed
fractions of that rate. Each steady-state run serves 400 requests and
is repeated thrice unless stated otherwise. Latency statistics discard
the first and last 10\% of requests and pool the remaining requests
across repetitions. We report fine-tuning throughput as completed
samples/s while requests are being served, together with tail
\ac{TTFT} and \ac{TPOT} relative to the corresponding serving-only
execution.

\paragraph{Fine-tuning workload.}
Every system trains the same job: LoRA rank 16 on the attention
projections, using alpaca-cleaned samples padded to 128 tokens and
AdamW with learning rate $2\cdot10^{-4}$. We adapt $(q,k,v,o)$ for
Qwen3-30B-A3B and Qwen1.5-MoE-A2.7B, and
$(q,kv_a,kv_b,o)$ for DeepSeek-V2-Lite.

\paragraph{Baselines.} We compare four systems.
\begin{itemize}
  \item \textbf{Serving only} provides the inference-latency reference.

  \item \textbf{\Uncoord} runs a stock torchtune LoRA trainer
  using FSDP alongside vLLM under CUDA MPS. Because its interference
  depends on the amount of GPU compute available to the trainer, we
  run it with the trainer's SM share capped at 6\% and 40\% and
  uncapped. The main text shows the 40\% cap;
  Appendix~\ref{sec:appx-naive} compares all three.

  \item \textbf{LLMStation}~\cite{llmstation} represents iteration-level colocation through suspendable tasklets. We update its implementation to the current vLLM (v1 engine) and PyTorch versions, and extend its training forward to Qwen3-MoE with shared serving weights. Its released fusion engine assumes dense projections with fixed weight placement, which does not hold for expert-parallel \ac{MoE}, especially with EPLB. We therefore use its separate PyTorch training path alongside inference under CUDA MPS with EPLB disabled. To obtain a representative baseline, we evaluated several expert-execution paths including a naive per-expert forward, reuse of \sys{}'s compiler-generated chunks and a packed-GEMM implementation. Executing padded batched GEMMs through cuBLAS gave the best LLMStation performance. This path does not use \sys{}'s compiler generated chunks or A100-specific grouped-GEMM kernel. LLMStation's latency ratios are measured relative to the same engine with training disabled, so differences between serving engines do not count against either system.

  \item \textbf{\sys{}} is the full system. It combines conventional
  decode/idle-time colocation with router-aware intra-layer bubble
  filling. Unless stated otherwise, it uses all mechanisms of
  \S\ref{sec:design}, including device-side readiness, paired chains,
  and early close.
\end{itemize}

We do not compare against colocation systems that use compute headroom
during prefill: prefill saturates the GPU at a few hundred
tokens~\cite{sarathi}, so at our prompt lengths that headroom is gone,
and the remaining serving-level opportunities, idle and decode time, are
the ones LLMStation targets.

\section{Additional experiments}
\label{sec:appx-experiments}

This section extends the evaluation of \S\ref{sec:evaluation} and
follows the order of its subsections.
For the end-to-end comparison (\S\ref{sec:eval-e2e}),
\S\ref{sec:appx-e2e} reports all three models and 1k--32k prompts,
\S\ref{sec:appx-tp2} repeats the Qwen3-30B-A3B comparison on two GPUs,
\S\ref{sec:appx-traces} describes both production traces and their full
results, and \S\ref{sec:appx-naive} sweeps the \ac{SM} cap of \uncoord.
For the sources of throughput and the latency protections
(\S\ref{sec:eval-ablation} and \S\ref{sec:eval-latency}),
\S\ref{sec:appx-ablation} repeats both breakdowns at other prompt
lengths, \S\ref{sec:appx-margin} selects the early-close margin,
\S\ref{sec:appx-threshold} tests the 4k-token bubble-filling threshold,
\S\ref{sec:appx-capacity} measures the decode and idle-time headroom
that serving leaves at each load, and \S\ref{sec:appx-bubble-use} shows
where the bubble capacity that \sys{} leaves unused goes.
\S\ref{sec:appx-eplb} extends the composition with \ac{EPLB}
(\S\ref{sec:eval-eplb}) to all workloads. Finally,
\S\ref{sec:eval-overheads} isolates the overhead of the \sys{} runtime
and verifies that chunked execution preserves training behavior.

\begin{figure*}[t]
  \centering
  \inputplot{ft_throughput}
  \caption{\textbf{Fine-tuning throughput while serving, all models.}
    Samples/s for \sys{}, LLMStation, and \uncoord (40\% cap)
    across three \ac{MoE} models, prompt lengths, and serving loads.
    Labels report \sys{}'s throughput relative to LLMStation.}
  \label{fig:appx-e2e-tput}
\end{figure*}

\begin{figure*}[t]
  \centering
  \inputplot{e2e_latency}
  \caption{\textbf{Serving latency under colocation, all models.}
    p90 \ac{TTFT} (top) and p90 \ac{TPOT} (bottom), relative to the
    corresponding serving-only execution. The dashed line marks serving
    alone and the dotted line a 10\% increase. Bars above 2$\times$ are
    clipped, with their value printed beside them. \Uncoord uses
    the 40\% cap.}
  \label{fig:appx-e2e-lat}
\end{figure*}

\subsection{End-to-end results for all models}
\label{sec:appx-e2e}

\S\ref{sec:eval-e2e} compares \sys{} with the baselines on Qwen3-30B-A3B
at 4k--32k prompts. Figures~\ref{fig:appx-e2e-tput} and~\ref{fig:appx-e2e-lat} extend
Figures~\ref{fig:eval-e2e-tput} and~\ref{fig:eval-e2e-lat} to
Qwen1.5-MoE-A2.7B and DeepSeek-V2-Lite, and to 1k, 2k, 8k, and 24k prompts.
Both models follow the same trends as Qwen3-30B-A3B: across all three models and 1k--32k prompts, \sys{} trains 1.9--6.2$\times$ faster than \llmstation{} (median 2.8$\times$). At 1k and 2k
tokens, prefills are too short for \sys{} to open useful windows,
so its gains there come from its runtime mechanisms on decode and
idle time (\S\ref{sec:eval-ablation}).

\sys{} keeps both p90 \ac{TTFT} and \ac{TPOT} within 10\% of serving
alone in 50 of the 63 workloads. Of the other 13, eight are the main
exception, \ac{TPOT} on Qwen1.5-MoE-A2.7B at long prompts: from 16k to
32k tokens, \sys{}'s p90 \ac{TPOT} reaches 1.14--1.47$\times$ serving
alone in every workload but 16k at medium load (1.09$\times$), highest
at 32k and medium load, while \ac{TTFT} stays within 6\%. Its short
decode iterations overlap with the fine-tuning worker's cross-GPU exchanges and
increase \ac{TPOT}. Three more exceed the bound by at most three
points: Qwen1.5-MoE-A2.7B at 8k and medium load (p90 \ac{TTFT}
1.13$\times$, \ac{TPOT} 1.12$\times$, where LLMStation's \ac{TPOT} is
1.15$\times$) and high load (\ac{TPOT} 1.10$\times$), and Qwen3-30B-A3B
at 2k and low load (\ac{TPOT} 1.11$\times$).

The last two are \ac{TTFT} at 1k-token prompts and high load:
Qwen1.5-MoE-A2.7B at 32 requests/s reaches 3.61$\times$ serving alone
and Qwen3-30B-A3B at 16 requests/s 1.13$\times$, while \ac{TPOT} stays
within 7\%. These prompts are below the 4k-token threshold, so \sys{}
trains only in decode and idle time, and the server is already near
saturation: on Qwen1.5-MoE-A2.7B each serving step takes about 29\,ms,
so 32 requests/s keep the GPU about 93\% busy, and serving alone
completes only 29 requests/s. Decode and idle-time training leaves the
median step unchanged but lengthens the slowest ones (p90 from 30 to
34\,ms in one run), and at this utilization that builds a request queue.
This is a limitation of the current policy: \sys{} does not pause
decode- and idle-time training while requests are queued, and we leave
such a queue-aware cutoff to future work.

\subsection{Smaller expert-parallel degree}
\label{sec:appx-tp2}

\S\ref{sec:eval-e2e} uses four GPUs; we test whether \sys{}'s advantage
holds with fewer expert-parallel GPUs, where each GPU hosts more
experts. Table~\ref{tab:appx-tp2} repeats the Qwen3-30B-A3B comparison with
half the GPUs, at \ac{TP}$=$\ac{EP}$=$2. High load is half the
\ac{TP}$=$\ac{EP}$=$4 high load of each prompt length, and medium load about a
quarter of it. \sys{} trains 3.3--3.9$\times$ faster than LLMStation
and keeps both p90 \ac{TTFT} and \ac{TPOT} within 10\% of serving alone
in every workload; the highest is p90 \ac{TPOT} at 8k tokens and high
load (1.10$\times$, against 1.11$\times$ for LLMStation).

\begin{table}[t]
  \centering
  \footnotesize
  \setlength{\tabcolsep}{1.6pt}
  \caption{\textbf{\ac{TP}$=$\ac{EP}$=$2.} Qwen3-30B-A3B, two
    repetitions. LLMS: LLMStation; \sysshort: \sys{}. FT: fine-tuning samples/s; TTFT and TPOT: p90 relative
    to serving alone.}
  \label{tab:appx-tp2}
  \begin{tabular}{llrrrrrrrr}
    \toprule
    & & & \multicolumn{2}{c}{FT} & & \multicolumn{2}{c}{LLMS} & \multicolumn{2}{c}{\sysshort} \\
    \cmidrule(lr){4-5}\cmidrule(lr){7-8}\cmidrule(lr){9-10}
    Len. & Load & req/s & LLMS & \sysshort & Ratio & TTFT & TPOT & TTFT & TPOT \\
    \midrule
    4k & med. & 2 & 1.58 & 5.60 & 3.5$\times$ & 1.05 & 1.08 & 1.04 & 1.05 \\
    4k & high & 4 & 0.95 & 3.46 & 3.6$\times$ & 1.02 & 1.01 & 1.02 & 1.00 \\
    8k & med. & 1 & 1.52 & 5.22 & 3.4$\times$ & 1.04 & 1.05 & 1.02 & 1.01 \\
    8k & high & 2 & 0.91 & 3.51 & 3.9$\times$ & 1.03 & 1.11 & 1.02 & 1.10 \\
    16k & med. & 0.5 & 1.46 & 4.86 & 3.3$\times$ & 1.02 & 1.02 & 1.02 & 1.04 \\
    16k & high & 0.875 & 0.90 & 3.18 & 3.5$\times$ & 1.02 & 1.02 & 1.01 & 1.01 \\
    \bottomrule
  \end{tabular}
\end{table}

\subsection{Production traces}
\label{sec:appx-traces}

\S\ref{sec:eval-e2e} summarizes \sys{}'s behavior under bursty production
arrivals, and \Cref{fig:intro-trace} shows one trace; here we describe
both traces and report full results. Constant-rate workloads let us
control prompt length and serving load, but production arrivals are
bursty. We therefore replay two traces from
the Azure LLM inference trace~\cite{splitwise}, \emph{code} and
\emph{conversation}, on Qwen3-30B-A3B, each over a 10-minute window
(minutes 18--28 and 17--27) at its original arrival rate. The traces
give each request's arrival time and its input and output lengths; we
fill prompts with arXiv text of that length. The code trace has 2,183
requests with long prompts and short outputs (median input 1,506 tokens,
p90 5,823; median output 13, p90 52) and very bursty arrivals: 3.9
requests/s on average but 0--14.6 per 10\,s interval, with a
coefficient of variation of inter-arrival times of 10.3. The
conversation trace has 3,716 requests at a steadier 6.2 requests/s
(3.9--8.2 per interval; coefficient of variation 1.0), with shorter
prompts (median 1,024, p90 4,080) and longer outputs (median 111, p90
421). Prompts of at least 4k tokens carry 47\% of prompt tokens in the
code trace and 10\% in the conversation trace.

\begin{figure}[t]
  \centering
  \inputplot{trace}
  \caption{\textbf{Trace-driven serving.}
    Request arrivals (10\,s bins, dashed: mean rate), fine-tuning
    throughput, and p90 serving latency relative to serving alone under
    two production traces. In the latency panels, the dashed line marks
    serving alone and the dotted line a 10\% increase. }
  \label{fig:eval-traces}
\end{figure}

Figure~\ref{fig:eval-traces} shows that \sys{} keeps its advantage
under bursty arrivals: it trains at 5.9 and 5.0 samples/s across the two
traces, compared with 1.8 and 1.5 samples/s for LLMStation, while p90
\ac{TTFT} and \ac{TPOT} remain within 7\% of serving alone.

\subsection{\Uncoord sweep}
\label{sec:appx-naive}

\S\ref{sec:eval-e2e} runs \uncoord with a 40\% \ac{SM} cap and notes that
lifting the cap does not help; here we justify that choice by sweeping
the cap. \Uncoord's interference depends on how much GPU compute its
trainer may use. Figures~\ref{fig:appx-naive-tput}
and~\ref{fig:appx-naive-lat} compare three settings of the trainer's
\ac{SM} cap under CUDA MPS: 6\%, 40\%, and uncapped. Raising the cap
from 6\% to 40\% raises the median fine-tuning throughput across the 63
workloads from 0.80 to 1.05 samples/s, but the number of workloads that
keep both p90 \ac{TTFT} and p90 \ac{TPOT} within 10\% of serving alone
falls from 51 of 63 to 22 of 63. Without a cap, median throughput does
not rise further (1.00 samples/s), and only 6 of 63 workloads stay
within that bound. In the median workload, \sys{} trains 5.9$\times$
faster than \uncoord at the 6\% cap and 4.4$\times$ faster than
at the 40\% cap. In a few high-load workloads \uncoord trains
faster than \sys{}, almost always by inflating p90 \ac{TTFT} by
5--31$\times$. The exception is DeepSeek-V2-Lite at 1k tokens and high
load, where \uncoord at the 40\% cap or uncapped trains about
1.2$\times$ faster than \sys{} while staying within 3\% of serving alone.

\begin{figure*}[t]
  \centering
  \inputplot{naive_sweep_tput}
  \caption{\textbf{\Uncoord sweep: fine-tuning throughput.}
    Samples/s of the torchtune trainer colocated with vLLM, with the
    trainer's \ac{SM} share capped at 6\%, at 40\%, or uncapped.}
  \label{fig:appx-naive-tput}
\end{figure*}

\begin{figure*}[t]
  \centering
  \inputplot{naive_sweep_lat}
  \caption{\textbf{\Uncoord sweep: serving latency.}
    p90 \ac{TTFT} (top) and p90 \ac{TPOT} (bottom) relative to serving
    alone, for the three \ac{SM} caps of
    Figure~\ref{fig:appx-naive-tput}. The dashed line marks serving alone
    and the dotted line a 10\% increase. Bars above 2$\times$ are clipped,
    with their value printed beside them.}
  \label{fig:appx-naive-lat}
\end{figure*}
\subsection{Ablation details}
\label{sec:appx-ablation}

\S\ref{sec:eval-ablation} and \S\ref{sec:eval-latency} break down \sys{}'s
throughput and its latency protections at 8k and 4k prompts; here we
check that both hold at other prompt lengths.
Figure~\ref{fig:appx-sources} extends Figure~\ref{fig:eval-sources} to
4k and 16k prompts. The pattern is the same at every length:
\cfgOptColo raises throughput over \cfgColo by 36--60\% at low and
medium load, bubble filling changes throughput by $-7$ to $-2$\% at
low load and adds 13--37\% at medium and high load. At 16k and high load
all three configurations train at most 0.53 samples/s, so the
relative gain of \cfgOptColo there (+159\%) is large but small in
absolute terms.

Figure~\ref{fig:appx-latency-ablation} extends
Figure~\ref{fig:eval-latency-ablation} to 8k and 16k prompts. Unthrottled
bubbles hurt serving only near saturation, and most at 16k prompts and
2 requests/s, where p90 \ac{TTFT} reaches 9.15$\times$ serving alone.
At 8k and 4 requests/s the server keeps enough headroom that the same
configuration keeps p90 \ac{TTFT} within 1.08$\times$, but its p90
\ac{TPOT} reaches 1.22$\times$. At 16k and high load, \sys{}
keeps \ac{TTFT} at 1.01$\times$ but its p90 \ac{TPOT} is 1.21$\times$.

\begin{figure*}[t]
  \begin{minipage}[t]{\columnwidth}
    \vspace{0pt}\centering
  \inputplot{latency_ablation}
  \caption{\textbf{Mechanisms that protect serving latency, 4k--16k
    prompts.} Rows: fine-tuning throughput, p90 \ac{TTFT} and p90
    \ac{TPOT} relative to serving alone. Configurations as in
    Figure~\ref{fig:eval-latency-ablation}.}
  \label{fig:appx-latency-ablation}
  \end{minipage}\hfill
  \begin{minipage}[t]{\columnwidth}
    \vspace{0pt}\centering
  \inputplot{sources}
  \caption{\textbf{Sources of \sys{}'s fine-tuning throughput, 4k--16k
    prompts.} Configurations as in Figure~\ref{fig:eval-sources}.}
  \label{fig:appx-sources}
  \end{minipage}
\end{figure*}

\subsection{Sensitivity to the early-close margin}
\label{sec:appx-margin}

In \S\ref{sec:eval-latency}, the early close is one of the protections
that keep training off the serving critical path; this subsection
selects its margin. \sys{} closes each admission window a margin $m$ before its predicted
end (\S\ref{sec:design-optC}). We select $m$ with the sweep in
Table~\ref{tab:appx-margin}, which varies $m$ from 0 to 120\,$\mu$s on
Qwen3-30B-A3B at three high-load workloads, where bubble filling
matters most (\S\ref{sec:eval-ablation}); all other settings are those
of \sys{}. The mechanism sets a floor on $m$: the margin should cover
one 50\,$\mu$s chunk, so that a chunk admitted just before
the close ends before the predicted end, and the predictor's p95 error
of 45\,$\mu$s (\S\ref{sec:design-predict}). The
sweep separates two regimes. Margins above 60\,$\mu$s cost 10--19\% of
fine-tuning throughput without lowering serving latency: p90 \ac{TTFT}
stays within 6\% of serving alone at every margin, and p90 \ac{TPOT}
does not fall (at 16k it rises to 1.22--1.25$\times$). Between 0 and
60\,$\mu$s, throughput and latency differ by no more than run-to-run
variation. Closing the window at the predicted end is what protects
serving (\S\ref{sec:eval-latency}); the margin absorbs prediction
error. We therefore pick $m=60\,\mu$s, the largest margin that costs no
measurable throughput: it clears the floor and gives the most headroom
against mispredictions at no cost.

\begin{table}[t]
  \centering
  \small
  \setlength{\tabcolsep}{4pt}
  \caption{\textbf{Early-close margin sweep.} Qwen3-30B-A3B,
    \ac{TP}$=$\ac{EP}$=$4. FT: fine-tuning samples/s; TPOT: p90
    \ac{TPOT} relative to serving alone. Two repetitions at 4k, one at
    8k and 16k; at 4k and 60\,$\mu$s, one run with a queueing episode
    (p90 \ac{TTFT} 559\,ms against 192\,ms) is excluded. Bold: the margin
    selected for \sys{}, 60\,$\mu$s.}
  \label{tab:appx-margin}
  \begin{tabular}{rrrrrrr}
    \toprule
    & \multicolumn{2}{c}{4k, 8 req/s} & \multicolumn{2}{c}{8k, 4 req/s} & \multicolumn{2}{c}{16k, 2 req/s} \\
    \cmidrule(lr){2-3}\cmidrule(lr){4-5}\cmidrule(lr){6-7}
    $m$ ($\mu$s) & FT & TPOT & FT & TPOT & FT & TPOT \\
    \midrule
    0 & 0.96 & 1.05 & 1.59 & 1.10 & 0.53 & 1.09 \\
    30 & 0.98 & 1.05 & 1.55 & 1.11 & 0.51 & 1.13 \\
    \textbf{60} & \textbf{0.97} & \textbf{1.05} & \textbf{1.63} & \textbf{1.08} & \textbf{0.51} & \textbf{1.08} \\
    90 & 0.86 & 1.06 & 1.39 & 1.13 & 0.43 & 1.22 \\
    120 & 0.83 & 1.05 & 1.47 & 1.09 & 0.42 & 1.25 \\
    \bottomrule
  \end{tabular}
\end{table}

\subsection{Bubble-filling threshold}
\label{sec:appx-threshold}

\cfgThrottled and \sys{} skip prefills shorter than 4k tokens
(\S\ref{sec:eval-latency}, \S\ref{sec:design-predict}); here we test
whether this threshold is needed. To do so, we rerun
Qwen3-30B-A3B at 2k-token prompts, where it never fires, with bubble
filling enabled in every prefill step. At low load (3 requests/s),
removing the threshold lowers fine-tuning throughput from 5.11 to 4.69
samples/s ($-8$\%), and p90 \ac{TTFT} and \ac{TPOT} rise from
1.07$\times$ and 1.05$\times$ serving alone to 1.13$\times$ and
1.09$\times$. At high load (12 requests/s), throughput falls from 1.57
to 1.39 samples/s ($-11$\%), and p90 \ac{TTFT} and \ac{TPOT} rise from
1.04$\times$ and 1.05$\times$ to 1.08$\times$ and 1.13$\times$. At 2k
tokens the bubbles are too short to amortize prediction and gating, so
filling them only adds contention. Both settings use two
repetitions, with latency relative to serving alone in the same
repetition.

\subsection{Spare capacity while serving}
\label{sec:appx-capacity}

\S\ref{sec:eval-ablation} argues that bubbles matter more as serving load
consumes conventional headroom; here we measure that headroom directly.
Figure~\ref{fig:appx-capacity} decomposes the GPU capacity a colocated
job could use while vLLM serves alone, for every model, prompt length,
and load from a quarter of the high load up to the high load. Idle time
is wall time with no forward step running. Decode headroom is the
fraction of the GPU a co-running process can use inside decode steps
without measurably slowing them, measured with a throttled GEMM probe
under CUDA MPS. Prefill bubbles are the per-GPU waits before each layer's
all-reduce, inside which \sys{} opens windows for fine-tuning. For prompts up to 16k tokens, decode
headroom accounts for 35--74\% of a GPU at the lightest load but only
3--11\% at high load, while bubble capacity grows with prefill activity.

\begin{figure*}[t]
  \centering
  \inputplot{capacity}
  \caption{\textbf{Spare GPU capacity while serving alone}, as a share of
    one GPU: idle time (whole GPU free), spare compute inside decode
    steps that a co-running process can use without slowing them, and
    prefill bubbles. vLLM alone, \ac{TP}$=$\ac{EP}$=4$, 400 requests per
    run; x labels give the request rate, from the high load down to a
    quarter of it for each prompt length.}
  \label{fig:appx-capacity}
\end{figure*}

\begin{figure*}[t]
  \begin{minipage}[t]{\columnwidth}
    \vspace{0pt}\centering
  \inputplot{bubble_supply}
  \caption{\textbf{Using expert-parallel bubbles.}
    Top: accounting of available non-straggler bubble time at
    8k tokens and 2 requests/s. Bottom: total bubble time created
    while different colocated execution policies run, normalized
    to serving alone, for \cfgColo (Coloc.), \cfgUnthr (Unthr.),
    \cfgThrottled (Thrott.), \sys{} (\sysshort), and \sys{} with one
    fine-tuning chain (\sysshort, 1 ch.); \S\ref{sec:eval-latency} defines
    the configurations. The accounting categories add up to 104\%: 46\,ms
    (3.8\% of bubble time) fall into two categories.}
  \label{fig:eval-bubble-use}
  \end{minipage}\hfill
  \begin{minipage}[t]{\columnwidth}
    \vspace{0pt}\centering
  \inputplot{eplb}
  \caption{\textbf{\sys{} composes with expert load balancing, all
    workloads.}
    Top: fine-tuning throughput without and with \ac{EPLB}
    (hatched), split into \cfgOptColo{} and what bubble filling adds;
    labels give the change from filling, and a dashed outline marks a
    loss. Middle and bottom: p90 \ac{TTFT} and p90 \ac{TPOT} of
    serving alone and \sys{}, relative to serving alone with \ac{EPLB}
    (dashed line; dotted: 10\% increase).}
  \label{fig:appx-eplb}
  \end{minipage}
\end{figure*}

\begin{figure*}[t]
  \centering
  \inputplot{sched_overhead}
  \caption{\textbf{Scheduling overhead with an idle fine-tuning worker.}
    Left three panels: p50 and p90 \ac{TTFT} and \ac{TPOT} with \sys{}'s
    hooks enabled but no training admitted, relative to serving alone,
    at 8k tokens and three loads per model (400 requests per run,
    pooled over runs). Right: per-layer cost of the prediction and
    admission hooks on one A100; solid bars are GPU time on the serving
    stream, hatched bars host time.}
  \label{fig:appx-sched-overhead}
\end{figure*}

\begin{figure*}[t]
  \centering
  \inputplot{overheads}
  \caption{\textbf{Memory overhead and training correctness.}
    Left: GPU memory taken per GPU, as a share of 80\,GB, by LLMStation
    (LLMS), \sys{} (\sysshort), \ac{EPLB}'s redundant experts, and both together
    (Both).
    vLLM has no \ac{EPLB} for Qwen1.5-MoE. Right: 25-step moving-average
    training loss on Qwen3-30B-A3B for \sys{} and a standalone HF+PEFT
    trainer with the same data, LoRA configuration, and learning rate.}
  \label{fig:appx-overheads}

\end{figure*}

\begin{figure}[t]
  \centering
  \inputplot{mem_breakdown}
  \caption{\textbf{Fine-tuning worker memory per GPU.} Breakdown of the resident
    \sys{} fine-tuning worker's footprint on Qwen3-30B-A3B and DeepSeek-V2-Lite.}
  \label{fig:appx-mem-breakdown}
\end{figure}

\subsection{Bubble utilization}
\label{sec:appx-bubble-use}

\S\ref{sec:eval-latency} notes that \sys{}'s protections leave some bubble
capacity unused; here we break down where that capacity goes. We define \emph{bubble utilization} as the fraction of measured
non-straggler bubble GPU-time during which fine-tuning computation
executes. Unlike the trace-driven experiment in Appendix~\ref{sec:appx-gap}, which
isolates what different planners could admit into recorded bubbles,
this measurement accounts for the complete \sys{} runtime: prediction,
admission, training dependencies, chunk readiness, and early close.

At 8k tokens and 2 requests/s, \sys{} executes useful fine-tuning work
during 36\% of the available bubble time. Approximately 50\% passes with
an admission window open but no fine-tuning chunk executing: the
fine-tuning worker is waiting on one of its own cross-GPU dependencies or has no ready
chunk. The remaining roughly 15\% is deliberately left unused because
\sys{} closes each window before its predicted end.

Open-but-idle time arises because \sys{} lets GPUs progress unevenly
(\S\ref{sec:design-compiler}) and does not try to equalize them.
Equalizing would rarely help: on Qwen3-30B-A3B at 16k tokens, the
most-loaded GPU changes between consecutive layers 70\% of the time,
and each of the four GPUs is the most-loaded in 14--39\% of layers,
so a GPU that falls behind in one layer often has a bubble to catch up in
the next, as well as during decode and idle time, which open on every
GPU. A GPU that runs ahead, however, can reach one of the training
step's cross-GPU exchanges before a lagging GPU has produced its
input, and its window then passes with nothing to run.

These losses have different causes. Open-but-idle time is a readiness
limitation: serving has exposed usable capacity, but training cannot
always supply work at the same fine granularity. Device-side
readiness and paired chains target this loss by keeping more work ready
(\S\ref{sec:design-chunk}), but they do not remove the dependencies
between fine-tuning chunks. Early-close time, in contrast, is capacity that
\sys{} gives up on purpose to avoid extending the serving critical path.

Colocated training could also create the capacity it later uses.
\cfgUnthr increases total bubble time by 17\% relative to serving
alone: training delays the most-loaded GPU, so its peers wait longer.
\sys{} increases it by only 4\% (2\% with one fine-tuning chain), because training
is admitted primarily on GPUs that would otherwise be waiting. The
windows also behave conservatively: 98\% close before the serving
dependency clears, with a median 57\,$\mu$s of remaining slack, and
the straggling GPU opens a window in only 4\% of layers.

\subsection{Composition with EPLB, all workloads}
\label{sec:appx-eplb}

\S\ref{sec:eval-eplb} shows that \sys{} composes with expert load
balancing at 8k prompts; here we report all six workloads.
Figure~\ref{fig:appx-eplb} extends Figure~\ref{fig:eval-eplb} to 4k, 8k,
and 16k prompts at medium and high load, with latency relative to
serving alone with \ac{EPLB}. The only workload outside 7\% of that
baseline is 16k prompts at 1 request/s, where p90 \ac{TPOT} reaches
1.32$\times$ (p90 \ac{TTFT} 1.06$\times$); there, bubble filling adds
no throughput ($-2$\%), so the residual bubbles are better left unused.
The tallest latency bar, 1.34$\times$ p90 \ac{TPOT} at 16k prompts and
2 requests/s, is serving alone \emph{without} \ac{EPLB}: it reflects
\ac{EPLB}'s own benefit to serving, not a cost of \sys{}, which stays
at 1.06$\times$ there.

\subsection{Overheads and correctness}
\label{sec:eval-overheads}

\S\ref{sec:evaluation} measures \sys{} with fine-tuning admitted; here we
isolate the cost of the runtime itself, independent of the fine-tuning
work it admits, and verify that its chunked execution preserves
training behavior.

\paragraph{Scheduling overhead.}
\sys{} places prediction and admission logic on the serving path, so
even a fine-tuning worker that never executes could in principle perturb
inference. To isolate this overhead, we run the complete runtime with
the fine-tuning worker resident but never admit a fine-tuning chunk. Each prefill
\ac{MoE} layer still predicts its bubble, opens an admission window,
arms early close, and closes the window.

Together, these hooks consume approximately 10\,$\mu$s of GPU time
per layer. Their host-side work takes approximately 130\,$\mu$s per
layer, but overlaps GPU execution; across all three models, measured
prefill-step duration changes by at most 0.5\%. Correspondingly, p50
and p90 \ac{TTFT} and \ac{TPOT} remain within 2.5\% of serving alone
in all but two configurations. On Qwen3-30B-A3B at medium load, p90
\ac{TPOT} rises by 6.8\% while p50 \ac{TPOT} and both \ac{TTFT}
percentiles stay within 1\%. The other exception is
DeepSeek-V2-Lite at its highest tested load, where the server is busy
approximately 95\% of the time and p50 and p90 \ac{TTFT} increase by
6.4\% and 7.0\%. Reducing the load from 7 to 6 requests/s removes this effect,
indicating that even sub-percent runtime overhead can amplify once the
server is already at the edge of queueing.
Figure~\ref{fig:appx-sched-overhead} reports the per-model results and
the per-layer hook cost.

\paragraph{Memory overhead.}
Keeping a separate fine-tuning worker resident reduces the memory available to the
serving KV cache by 4.3--5.8\,GiB per GPU, or 5--7\% of an 80\,GB A100,
compared with vLLM alone. The fine-tuning worker itself holds 3.0--4.6\,GiB
of this (Figure~\ref{fig:appx-mem-breakdown}); the rest is buffers,
CUDA contexts, and exchange mailboxes outside its allocator. The LoRA
parameters, gradients, and optimizer state take at most 0.15\,GiB;
most of the fine-tuning worker's memory is the activation pool backing the
compiled training graphs and its separate CUDA context. For comparison,
LLMStation reduces the KV cache by 1.5--2.4\,GiB per GPU because its
trainer executes inside the serving process and can share the CUDA
context and allocator. Figure~\ref{fig:appx-overheads} (left) compares
these reductions, including the additional cost when \ac{EPLB} is
enabled.

\paragraph{Training correctness and runtime reliability.}
\label{appendix:correctness}
Chunking a training step across short windows and two fine-tuning chains
must not change the optimization being performed. We therefore train
Qwen3-30B-A3B with \sys{} and with a standalone HF+PEFT trainer using
the same data, LoRA configuration, and learning rate. Across 476
steps, the 25-step moving-average losses differ by 1.2\% on average
and at most 4.5\%. We additionally observed 239,218 device-side
exchanges across 320 \sys{} runs with no mismatch, fallback, or
timeout (Figure~\ref{fig:appx-overheads}, right). Thus, the fine-grained execution required to exploit short
bubbles preserves the expected training behavior while
operating reliably across long-running experiments.

